\documentclass[10pt,a4paper]{amsart}
\usepackage{amsmath,amssymb,amsthm,graphicx}
\usepackage{amsmath}
\usepackage{amssymb}
\usepackage{hyperref}
\newtheorem{Theorem}{Theorem}[section]
\newtheorem{Corollary}[Theorem]{Corollary}
\newtheorem{Lemma}[Theorem]{Lemma}
\newtheorem{Proposition}[Theorem]{Proposition}
{\theoremstyle{definition}
\newtheorem{Note}[Theorem]{Note}
\newtheorem{Definition}[Theorem]{Definition}
\newtheorem{Problem}[Theorem]{Problem}
}
\usepackage{cleveref}
\crefname{originaldefinition}{original definition}{original definitions}
\Crefname{originaldefinition}{Original definition}{Original definitions}
\makeatletter
\newlabel{def:uaw}{{1.2}{source}{}{sourcepointer-1}{}}
\newlabel{def:uaw@cref}{{[originaldefinition][0][]1.2}{[1][0][]source}{}{}{}}
\makeatother
\title{The universal Askey-Wilson algebra maps onto the polynomial algebra in three variables}
\author{Paul Terwilliger}
\date{doi:10.3842/SIGMA.2011.069; CC BY4.0}
\begin{document}
\maketitle
\noindent\textit{Editorial note.} All delivered mathematical text is the author's own, copied from the pinned publisher source sources/sigma-2011-069-askey-wilson/pinned-original-source.tex; no statement or proof step was rewritten, reordered or removed.\par\medskip
\noindent\textit{Editorial note.} Preamble provenance: the retained preamble is the source's own layout adapter plus the author's macro block at lines 10-18 of sources/sigma-2011-069-askey-wilson/pinned-original-source.tex (everything between the publisher class metadata and \textbackslash{}begin\{document\}, comments excluded). Nothing was added, no macro was invented, and the publisher document class is deliberately not reproduced; the excerpt serves only as preamble, never as a span coordinate.\par\medskip
\noindent\textit{Editorial note.} External references are navigation pointers only; every deliberately kept step of the author's argument stays inside the counted spans. The displayed numbers of those pointers are the author's own numbering under the author's \textbackslash{}newtheorem declarations (the source has no subsections, so numbering is section.index), computed mechanically by the lane's scanner.\par\medskip
\noindent\textit{Editorial note.} The single reference to another labelled item is Definition def:uaw, recorded as a navigation stub of kind definition carrying the author's own number; the definition itself is not reproduced, and the step that consumes it stays inside the counted span.\par\medskip
\noindent\textit{Editorial note.} The counted spans contain no bibliography citation, so no bibliography excerpt is attached to this card.\par\medskip
\noindent\textit{Editorial note.} Licence binding: the card's source\_version now carries the value recorded in this source's own licence evidence (sources/sigma-2011-069-askey-wilson/license-observation.json), namely 'SIGMA7(2011),069; publisher final form2011-07-09'. The adopt step of the publication line compares the two and refused the earlier revision; the card is aligned to the evidence rather than the evidence to the card.\par\medskip
\noindent\textit{Allowed standard prerequisites (editorial).}\begin{itemize}
\item The defining relations of the universal Askey-Wilson algebra Delta. Same source, Definition labelled def:uaw; invoked inside the counted proof at lines 2711-2713. Supplies the relations that the three barred generators must satisfy for the homomorphism to exist.
\item The generators A, B, C of Delta. Same source, the algebra Delta is generated by A, B, C and the barred images generate the polynomial algebra; used at lines 2714-2731. Gives uniqueness (A, B, C generate the source) and surjectivity (their images generate the target).
\end{itemize}\medskip
\ifdefined\hypertarget\hypertarget{sourcepointer-1}{}\fi
\noindent\textit{Source reference.} Original definition 1.2: sources/sigma-2011-069-askey-wilson/pinned-original-source.tex, source lines 2713--2713. Navigation pointer for \textbackslash{}ref\{def:uaw\} used inside the counted span (line 2713); the argument that consumes it is retained in the span.\par\medskip
\begin{Lemma}
\label{lem:com2}
There exists a unique $\mathbb F$-algebra homomorphism
$\Delta \to
\mathbb F\lbrack
\overline A,
\overline B,
\overline C
\rbrack $
that sends
\begin{gather*}
A \mapsto \overline A,
\qquad
B \mapsto \overline B,
\qquad
C \mapsto \overline C.
\end{gather*}
This homomorphism is surjective.
\end{Lemma}

\begin{proof}
The algebra
$
\mathbb F\lbrack
\overline A,
\overline B,
\overline C
\rbrack $ is commutative, so each of
\begin{gather*}
 \overline A+ \frac{q\overline B \,
\overline C-q^{-1}\overline C \, \overline B}{q^2-q^{-2}},
\qquad
\overline B+
\frac{q\overline C\, \overline A-q^{-1}\overline A \,\overline C}{q^2-q^{-2}},
\qquad
\overline C+
\frac{q\overline A\,\overline B-q^{-1}\overline B\,\overline A}{q^2-q^{-2}}
\end{gather*}
is central in
$
\mathbb F\lbrack
\overline A,
\overline B,
\overline C
\rbrack $.
Consequently
$\overline A$,
$\overline B$,
$\overline C
$ satisfy the def\/ining relations for~$\Delta $ from
Def\/inition~\ref{def:uaw}. Therefore the homomorphism exists.
The homomorphism is unique since
$A$,
$B$,
$C$ generate~$\Delta $.
The homomorphism is surjective
since
$
\overline A$,
$\overline B$,
$\overline C
$
generate
$
\mathbb F\lbrack
\overline A,
\overline B,
\overline C
\rbrack $.
\end{proof}

\begin{thebibliography}{99}

\footnotesize\itemsep=0pt
\bibitem{alpern}
Alperin R.C.,
Notes: ${\rm  {PSL}}_2(\mathbb Z)= \mathbb Z_2 \star \mathbb Z_3$,
\href{http://dx.doi.org/10.2307/2324963}{{\it Amer. Math. Monthly}} {\bf 100} (1993), 385--386.
\end{thebibliography}
\end{document}
